\section*{Original Appendix A: Missed Lemmas (pp.17:21--25)}
\subsection*{A.1 Lemma 18}
At first we prove an auxiliary Lemma.
\begin{result}{Lemma 28}
If $G:=(V,E)$ satisfies $(r,a)$-expansion property, $M\subseteq E$, and $S\subseteq V$ of size at most $r$, such that $|E(S,\overline S)\setminus M|\le b\Delta|S|$ then $|S|\le\frac{|M|}{(a-b)\Delta}$.
\end{result}
\begin{proof}
The size of $S$ is at most $r$, hence:
\[b\Delta|S|\ge|E(S,\overline S)\setminus M|\ge a\Delta|S|-|M|.\]
Thus $|S|\le\frac{|M|}{(a-b)\Delta}$.
\end{proof}
\begin{result}{Lemma A.1 (18)}
Let $(p,U_p,\sigma_p)$ and $(q,U_q,\sigma_q)$ be $0.9$-good triples. Then there is an assignment $\kappa$ such that:
\begin{itemize}\item for any $\gamma$ that is an instance of $\sigma_p\joint\sigma_q$ an assignment $\gamma\cup\kappa$ is a $\beta'$-mlce on $\supp(\sigma_p)\cup\supp(\sigma_q)$;
\item $|\supp(\gamma\cup\kappa)|\le\nu_{k-1}\Delta r$.\end{itemize}
Moreover if $p=q$ then triple $(p,U_p,\sigma_p\cup\kappa)$ is $\beta'$-good.
\end{result}
\begin{proof}
Let $S:=V_{\sigma_p}\cup V_{\sigma_q}$, $E_\sigma:=\supp(\sigma_p)\cup\supp(\sigma_q)$ and $B\subseteq V\setminus S$ be the maximal set of vertices that satisfies:
\[|B|\le r;\qquad |E(B,\overline B)\setminus E_\sigma|\le\beta'\Delta|B|.\]
At first we give an upper bound on the size of set $B$.
Partial assignment $\sigma_p$ is $0.9$-mlce on $M\cup U_p$. $\beta'\Delta|V_{\sigma_p}|\ge|E(V_{\sigma_p},\overline{V_{\sigma_p}})\setminus(M\cup U_p)|$ and by Lemma 28
\[|V_{\sigma_p}|\le\frac{|M\cup U_p|}{(\beta-\beta')\Delta}\le2\frac{\nu_k}{\beta-\beta'}r\le\tfrac12\nu_{k-1}r.\]
By analogy the same holds for $V_{\sigma_q}$.
The equality $E(B,\overline B)\cap E_\sigma=E(B,S)\cup(E(B)\cap|M\cup U_p\cup U_q|)$ together with $|E(B,\overline B)\setminus E_\sigma|\le\beta'\Delta|B|$ implies:
\[(1-\beta')\Delta|B|-|M\cup U_p\cup U_q|\le|E(B,S)|.\]
By Mixing Lemma:
\[|E(B,S)|\le\frac\Delta n|B||S|+\alpha\Delta\sqrt{|S||B|}.\]
For the sake of contradiction assume that $|B|\ge|S|$ thus:
\[|E(B,S)|\le\frac\Delta n|B||S|+\alpha\Delta|B|.\]
Altogether:
\[(1-\beta')\Delta|B|\le\frac{\Delta r}{n}\nu_{k-1}|B|+\alpha\Delta|B|+3\nu_k\Delta r\le2\alpha\Delta|B|,\]
that contradicts the choice of $\alpha$ and $\beta'$, hence $|B|\le|S|\le\nu_{k-1}r$.
At first we show that $(V\setminus(S\cup B),E\setminus(E_\sigma\cup E(B)))$ satisfies $(r,\beta')$-expansion property. By contradiction, suppose that there is a set $B'\subseteq V\setminus(S\cup B)$ of size at most $r$ such that
\[|E(B',\overline{B'})\setminus(E_\sigma\cup E(B))|<\beta'\Delta|B'|.\]
Again by Lemma 28 we conclude that:
\[|B'|\le\frac{|M\cup U_p\cup U_q|+\Delta|S\cup B|}{(\beta-\beta')\Delta}\le\nu_{k-1}r+\tfrac12r\le\tfrac34r.\]
But it implies that $|B\cup B'|\le r$, moreover:
\begin{align*}
|E(B\cup B',\overline{B\cup B'})\setminus E_\sigma|&\le\beta'\Delta|B|+\beta'\Delta|B'|\\
&=\beta'\Delta|B\cup B'|.\qquad\text{$B$ and $B'$ are disjoint}
\end{align*}
That contradicts the choice of $B$.
Now we find a proper assignment $\kappa$ on the $E(B)\setminus E_\sigma$ such that for all $v\in B$:
\[\sum_{e\in E:\st(e)=v}x_e\ge0.8\cdot\Delta.\]
Since $\sigma_p$ is an $(r,0.6,0.9)$-locally consistent assignment, then $(V\setminus V_{\sigma_p},E\setminus\supp(\sigma_p))$ satisfies $(r,0.9)$-expansion property. By analogy we have the same property for $\sigma_q$ that implies: $(V\setminus S,E\setminus E_\sigma)$ satisfies $(r,0.8)$-expansion property. Indeed, consider a set $C\subseteq V\setminus S$ of size at most $r$:
\begin{align*}
|E(C,\overline C)\setminus E_\sigma|&=|E(C,\overline C)|-|E(C,\overline C)\cap E_\sigma|\\
&\ge|E(C,\overline C)|-|E(C,\overline C)\cap\supp(\sigma_p)|-|E(C,\overline C)\cap\supp(\sigma_q)|\\
&=|E(C,\overline C)\setminus\supp(\sigma_p)|-0.1\cdot\Delta|C|\\
&\ge0.8\cdot\Delta|C|.
\end{align*}
By Proposition 12 there is an enumeration of vertices in $B$: $v_1,v_2,\ldots,v_{|B|}\in B$ and a sequence $R_1,\ldots,R_{|B|}\subseteq(E(B)\setminus E_\sigma)$ such that:
\[R_i=E(\{v_i\},V\setminus\{v_1,v_2,\ldots,v_i\})\setminus E_\sigma;\quad |R_i|\ge0.8\Delta.\]
We define $\kappa$ in the following way:
\begin{itemize}\item for an $e\in R_i$ we assign corresponding variables to direct the flow outside of the vertex $v_i$ (i.e. if $e'$ is a directed copy of $e$ that goes outside of $v_i$ we set $x_{e'}$ to $1$ and set the dual edge to $0$);
\item for all loops inside the set $B$ we assign corresponding variables to $0$.\end{itemize}
Let $\gamma$ be an instance of $\sigma_p\joint\sigma_q$, $\zeta:=\gamma\cup\kappa$ and $V_\zeta:=S\cup B$. We have already shown that the graph $(V\setminus V_\zeta,E\setminus\supp(\zeta))$ satisfies $(r,\beta')$-expansion property. We want to show that vertices in $V_\zeta$ are $0.6$-satisfied by $\zeta$. Consider four cases.
\begin{enumerate}
\item $v\in V_\rho$. Both assignments $\sigma_p$ and $\sigma_q$ extend an assignment $\rho$ hence $\gamma$ agreed with both assignments on edges incident to $V_\rho$. Thus $\gamma$ $0.6$-satisfies $v$.
\item $v\in V_{\sigma_p}\setminus V_\rho$. Let $E_v$ be a set of edges that are incident to $v$. At least $0.8\cdot\Delta$ of those edges carry outgoing flow from $v$ in $\sigma_p$. Denote those edges as $E_{\sigma_p}$.
If $v\notin V_{\sigma_q}$ then $\sigma_q$ may assign at most $0.1\cdot\Delta$ edges in $E_v$. That means that in $\gamma$ at least $0.7\cdot\Delta$ edges from $E_{\sigma_p}$ still carry outgoing flow from $v$.
If $v\in V_{\sigma_q}$ then $\sigma_p$ and $\sigma_q$ both $0.8$-satisfy $v$. Let $E_{\sigma_q}\subseteq E_v$ be the set of edges that carry outgoing flow from $v$ in $\sigma_q$. Then $E_{\sigma_p}\cap E_{\sigma_q}\ge0.6\cdot\Delta$, and all those edges carry outgoing flow from $v$ in $\gamma$.
Note that if $\sigma_p=\sigma_q$, then we $0.8$-satisfy $v$.
\item $v\in V_{\sigma_p}\setminus V_\rho$. By analogy with the previous case.
\item $v\in B$. We direct the flow on at least $0.8\cdot\Delta$ edges from $E_v$ outside of $v$ hence $\kappa$ $0.8$-satisfies $v$.
\end{enumerate}
By construction $V_\zeta:=V_{\sigma_p}\cup V_{\sigma_q}\cup B$ hence $|V_\zeta|\le\nu_{k-1}r$ and $|\supp(\zeta)|\le\nu_{k-1}r$. In order to check that $\zeta$ is $\beta'$-mlce note that:
\[|E(B,\overline B)\setminus E_\sigma|\le\beta'\Delta|B|\le\beta'\Delta|V_\zeta|,\]
but
\[|E(B,\overline B)\setminus E_\sigma|=|E(V_\zeta,\overline{V_\zeta})\setminus E_\sigma|\]
since $\sigma_p$ and $\sigma_q$ together assign all edges that are incident to $V_{\sigma_p}\cup V_{\sigma_q}$. Thus:
\[|E(V_\zeta,\overline{V_\zeta})\setminus E_\sigma|\le\beta'\Delta|V_\zeta|\]
that concludes the proof.
In case of $(p,U_p,\sigma_p)=(q,U_q,\sigma_q)$ it remains to show that $\zeta$ is $\beta'$-mlce on $M\cup U_p$. Again we note that:
\[|E(B,\overline B)\setminus E_\sigma|\le\beta'\Delta|B|,\]
and also
\[|E(V_{\sigma_p},\overline{V_{\sigma_p}})\setminus E_\sigma|\le\beta'\Delta|V_{\sigma_p}|,\]
hence
\begin{align*}
|E(V_{\sigma_p}\cup B,\overline{V_{\sigma_p}\cup B})\setminus E_\sigma|
&\le\beta'\Delta(|V_{\sigma_p}|+|B|)\\
&\le\beta'\Delta|V_{\sigma_p}\cup B|,
\end{align*}
where the last inequality holds since $B$ and $V_{\sigma_p}$ are disjoint, that concludes the proof.
\end{proof}
\subsection*{A.2 Lemma 23}
\begin{result}{Lemma A.2 (23)}
For all $i\le\ell$:
\begin{itemize}\item $\kappa_i$ exists;
\item $|V_{\tau_i}|\le\frac{1}{(\beta-\beta')\Delta}(\supp(\zeta)+\Delta|U_i|)$ and hence $|\tau_i|\le\frac{2}{\beta-\beta'}(|\supp(\zeta)|+\Delta|U_i|)$;
\item $(V\setminus V_{\tau_i},E\setminus\supp(\tau_i))$ satisfies $(r,\beta')$-expansion property.\end{itemize}
\end{result}
\begin{proof}
We show by induction on $i$ that:
\begin{itemize}\item $(V\setminus V_{\tau_i},E\setminus\supp(\tau_i))$ satisfies $(r,\beta')$-expansion property;
\item $|E(V_{\tau_i},\overline{V_{\tau_i}})\setminus(\supp(\zeta)\cup E(U_i))|<\beta'\Delta|V_{\tau_i}|$;
\item $|V_{\tau_i}|\le\frac{1}{(\beta-\beta')\Delta}(\supp(\zeta)+\Delta|U_i|)$ and hence $|\tau_i|\le\frac{2}{\beta-\beta'}(|\supp(\zeta)|+\Delta|U_i|)$.\end{itemize}
Assignment $\tau_0$ is $\zeta$ and $\zeta$ is $(r,0.6,\beta')$-locally consistent, in particular, $(V\setminus V_\zeta,E\setminus\supp(\zeta))$ satisfies $(r,\beta')$-expansion property and $E(V_\zeta,\overline{V_\zeta})\setminus\supp(\zeta))=\emptyset$.
By definition of $H_i$:
\[\beta'\Delta|H_i|>|E(H_i,\overline{H_i}\setminus\{u_i\})\setminus\supp(\tau_i)|\ge|E(H_i,\overline{H_i})\setminus(\supp(\tau_i)\cup E(H_i,\{u_i\}))|\]
and by Lemma 28
\[|H_i|\le\frac{|\supp(\tau_i)\cup E(H_i,u_i)|}{(\beta-\beta')\Delta}\le\frac{|\supp(\tau_i)\cup E(H_i,u_i)|}{(\beta-\beta')\Delta}\le\tfrac12\nu_{k-2}(r+1).\]
Hence $|H_i\cup V_{\tau_i}\cup\{u_i\}|\le r$ that together with:
\begin{align*}
&|E(H_i\cup V_{\tau_i},\overline{H_i\cup V_{\tau_i}})\setminus(\supp(\zeta)\cup E(U_{i+1}))|\\
&\le |E(V_{\tau_i},\overline{H_i\cup V_{\tau_i}})\setminus(\supp(\zeta)\cup E(U_{i+1}))|\\
&\quad+|E(H_i,\overline{H_i})\setminus(\supp(\zeta)\cup E(U_{i+1})\cup E(V_{\tau_i})|\\
&\le\beta'\Delta|V_{\tau_i}|+\beta'\Delta|H_i|\le\beta'\Delta|H_i\cup V_{\tau_i}|\le\beta'\Delta|H_i\cup V_{\tau_i}\cup\{u_i\}|
\end{align*}
implies $|V_{\tau_{i+1}}|=|H_i\cup V_{\tau_i}\cup\{u_i\}|\le\frac{1}{(\beta-\beta')\Delta}(|\supp(\zeta)|+\Delta|U_{i+1}|)$ by Lemma 28. Also $|\tau_{i+1}|\le\frac{2}{\beta-\beta'}(|\supp(\zeta)|+\Delta|U_{i+1}|)$ since by construction $\tau_{i+1}$ assigns only edges in $\supp(\zeta)\cup E(U_i\cup V_{\tau_{i+1}})$.
Now we show that a graph $(V\setminus V_{\tau_{i+1}},E\setminus\supp(\tau_{i+1}))$ satisfies $(r,\beta')$-expansion property. For the sake of contradiction assume that there is a set $S\subseteq V\setminus V_{\tau_{i+1}}$ of size at most $r$ such that: $E(S,\overline S)\setminus\supp(\tau_{i+1})\le\beta'\Delta|B|$.
By Lemma 28 $|S|\le\frac{|\supp(\tau_{i+1})|}{(\beta-\beta')\Delta}\le\tfrac12\nu_{k-2}(r+1)$. Hence $|H_i\cup S|\le r$ that together with:
\begin{align*}
&E(H_i\cup S,\overline{H_i\cup S}\setminus\{u_i\})\setminus\supp(\tau_i)|\\
&\le E(H_i,\overline{H_i\cup S}\setminus\{u_i\})\setminus\supp(\tau_i)|+E(S,\overline{H_i\cup S}\setminus\{u_i\})\setminus\supp(\tau_i)|\\
&\le\beta'\Delta|H_i|+\beta'\Delta|S|=\beta'\Delta|H_i\cup S|
\end{align*}
contradicts the choice of $H_i$.
To conclude the proof we have to show the existence of $\kappa_i$. Note that $(V\setminus V_{\tau_i},E\setminus\supp(\tau_i))$ satisfies $(r,\beta')$-expansion property. Consider an arbitrary set $B\subseteq V\setminus(V_{\tau_i}\cup\{u_i\})$ of size at most $r$:
\[|E(B,\overline B)\setminus(\supp(\tau_i)\cup E(\{u_i\}))|\ge\beta'\Delta|B|-E(B,\{u\}).\]
By Mixing Lemma:
\[|E(B,\{u\})|\le\frac\Delta n|B|+\alpha\Delta\sqrt B\le0.05\cdot\Delta|B|,\]
and hence
\[|E(B,\overline B)\setminus(\supp(\tau_i)\cup E(\{u_i\}))|\ge0.9\cdot\Delta|B|\]
and graph $(V\setminus V_{\tau_i}\setminus\{u_i\},E\setminus\supp(\tau_i))$ satisfies $(r,0.9)$-expansion property.
By Proposition 12 there is an enumeration of vertices in $H_i$: $v_1,v_2,\ldots,v_{|H_i|}\in H_i$ and a sequence $R_1,\ldots,R_{|H_i|}\subseteq E(H_i)\setminus(\supp(\tau_i)\cup E(\{u_i\}))$ such that:
\[R_k=E(\{v_k\},V\setminus\{v_1,v_2,\ldots,v_k\})\setminus(\supp(\tau_i)\cup E(\{u_i\}));\qquad |R_i|\ge0.9\cdot\Delta.\]
We define $\kappa_i$ for vertices $v_1,\ldots,v_{H_i}$ step by step, such that $\kappa_i$ on $E(v_k)$ satisfies the constraint:
\[\sum_{e\in E:\st(e)=v_k}(\tau'\cup\kappa_i)(x_e)-\sum_{e\in E:\en(e)=v_k}(\tau'\cup\kappa_i)(x_e)=c(v_k).\]
Since we have an access to the $0.9\cdot\Delta$ edges and others are already assigned, we can always choose the right values (loops are always assigned to zero).
\end{proof}
\section*{Original Appendix B: Garland in the Paths (p.17:25)}
\begin{result}{Lemma B.1 (19)}
There are $(p,U_p,\sigma_p),(q,U_q,\sigma_q)\in\mathcal S$ such that $(p,q)$ forms a $(k+1)$-garland.
\end{result}
\begin{proof}
Note that we can describe elements in $\mathcal P$ by a sequence of bits of size $s:=\nu_k\Delta r$. Each bit of this sequence describes an assignment for an edge $e$ that we choose on ``branching step''. From the construction it follows that different sequences generate different paths in the branching program and hence different elements of $\mathcal P$.
Let $s_k:=\lfloor s/(k+1)\rfloor$. We construct our garland by the iterative algorithm. After $i$-th iteration we have a set $S_i$ of sequences of size $is_k$ such that any two of the corresponding paths form $i$-garland and all paths end in the same node. The size of $S_i$ will be at least $\exp[s_k-\frac{i}{2k}s_k]$ for all $1\le i\le k+1$.
\begin{enumerate}
\item For $i=1$ consider all possible strings of length $s_k$ and paths that correspond to them. The branching program has size at most $2^{s_k/(2k)}$, hence there exists a node such that at least $2^{s_k(2k-1)/(2k)}$ paths end there. The set $S_1$ consists of all corresponding sequences.
\item For the step $i$, $2\le i\le k+1$, we consider all sequences in $S_{i-1}$. Let $v$ be the end node of all paths corresponding to sequences in the set. To each sequence $s\in S_{i-1}$ we append a string $u_s$ of $s_k$ bits in such a way that for any pair $r,r'\in S_{i-1}$ paths that corresponds to $ru_r$ and $r'u_{r'}$ differ at some node after $v$. Since $2^{s_k}\ge|S_{i-1}|$, it is possible to do this.
For the resulting sequences, we consider the set of the corresponding paths. The set of paths has size at least $2^{s_k(2k-i+1)/(2k)}$, and the size of the program is at most $2^{s_k/(2k)}$. Hence there exists a node such that $2^{s_k(2k-i)/(2k)}$ paths end there. Let $S_i$ be the set of sequences corresponding to those paths.
\end{enumerate}
After $k+1$ steps we have a set $S_{k+1}$, $|S_{k+1}|\ge2$, such that any two sequences in it correspond to a $(k+1)$-garland.
\end{proof}
